\section{Event detection}\label{sec:event-detection}

Basically, detecting that a rule application creates (or splits, merges, …) an orbit implies comparing the graph before and after this application~\cite{tierney2017using}.
%
However, comparing the members of a rule against each other allows a static detection of topological events.
Such a method presents the advantage of being independent of the object actually transformed by the rule.

\subsection{Creation}

\label{sec:creation}

A rule node matches at least one dart.
So when a rule creates a node, it also creates the corresponding darts.
If the right-hand side of a rule contains a created orbit, then applying this rule on any object creates one or several orbits.
For example, the triangulation rule (Fig.~\ref{fig:rule-triangulation}) creates the edge orbit $R\orbit{\alpha 0, \alpha 2}(n_{1})$.
Therefore, the application of this rule on the square face (Fig.~\ref{fig:before-after-house-triangulation}) creates four edge orbits $H\orbit{\alpha 0, \alpha 2}(a_{1})$, $H\orbit{\alpha 0, \alpha 2}(c_{1})$, $H\orbit{\alpha 0, \alpha 2}(e_{1})$ and $H\orbit{\alpha 0, \alpha 2}(g_{1})$ in $H$.
\begin{definition}[Orbit creation]
  \label{def:orbit-creation}
  Let $r : L \longrightarrow R$ be a rule, $m : L \to G$ a match, $t: G \longrightarrow^{r,m} H$ the transformation of $G$ by $(r,m)$, $\orbit{o}$ an orbit type, $n$ a created node of $R$ and $d$ a dart of $H$ matched by $n$.

  An orbit $R\orbit{o}(n)$ is \emph{created} in the rule $r$ if all of its nodes are created in $r$.

  An orbit $H\orbit{o}(d)$ is \emph{created} in the transformation $t$ if and only if all of its darts are created by $t$.
\end{definition}

\begin{proposition}
Let $r : L \longrightarrow R$ be a rule, $m : L \to G$ a match, $t: G \longrightarrow^{r,m} H$ the transformation of $G$ by $(r,m)$, $\orbit{o}$ an orbit type, $n$ a created node of $R$ and $d$ a dart of $H$ matched by $n$.

The orbit $H\orbit{o}(d)$ is created if and only if $R\orbit{o}(n)$ is created.
\end{proposition}

\begin{proofidea}
    Thanks to the syntactic conditions on rules, the orbit $R\orbit{o}(n)$ is complete.
    Thus, due to application of rules, $H\orbit{o}(d)$ is included in the image of $R\orbit{o}(n)$.
    Therefore, all darts of $H\orbit{o}(d)$ are created.
    Conversely, if $H\orbit{o}(d)$ is created, then this orbit is included in the image of $R\orbit{o}(n)$ and $R\orbit{o}(n)$ is created.
\end{proofidea}

Back to the example of face triangulation (Fig.~\ref{fig:rule-triangulation}~and~\ref{fig:before-after-house-triangulation}),
we follow the creation of an edge.
The edge orbit $R\orbit{\alpha 0, \alpha 2}(n_{1})$ in the triangulation rule is created because both its nodes $n_{1}$ and $n_{2}$ are created.
In the transformation, $n_{1}$ matches eight darts, including for example dart $a_{1}$, and $n_{2}$ also matches eight darts including dart $a_{2}$.
The edge orbit $H\orbit{\alpha 0, \alpha 2}(a_{1})$ is created as well as darts $a_{2}, h_{1}$ and $h_{2}$
(Fig.~\ref{fig:house-triangulation-right-match}).


\subsection{Split}

\begin{figure}[t]
    \centering
    \maxsizebox{\columnwidth}{.15\columnwidth}{\includegraphics{rule-unsewA2.pdf}}
    \caption{Rule $r : L \longrightarrow R$ of $\alpha 2$ unsewing along an edge ($\orbit{\alpha 0, \alpha 2}$-orbit)}
    \label{fig:rule-unsew-a2}
\end{figure}

\begin{figure*}[t]
  \begin{subfigure}[t]{.3\textwidth}
    \centering
      \includegraphics[height=.2\textheight]{demo-fig-square-plan-edge.pdf}
    \caption{Graph $G$}
    \label{fig:unsew-a2-left-match}
  \end{subfigure}
  \hfill
  \begin{subfigure}[t]{.3\textwidth}
    \centering
      \includegraphics[width=.2\textheight]{demo-fig-square-plan-edge-split.pdf}
    \caption{Graph $H_{1}$}
    \label{fig:unsew-a2-first-unsew}
  \end{subfigure}
  \hfill
  \begin{subfigure}[t]{.3\textwidth}
    \centering
      \includegraphics[height=.2\textheight]{demo-fig-square-plan-edge-split-2.pdf}
    \caption{Graph $H_{2}$}
    \label{fig:unsew-a2-second-unsew}
  \end{subfigure}
  \caption{Two successive applications of the $\alpha 2$ unsewing rule
    (Fig.~\ref{fig:rule-unsew-a2}) on a surface}
  \label{fig:before-after-unsew-a2}
\end{figure*}

When a rule explicitly splits an orbit, it also splits the set of matched darts.
%Such a split can occur either explicitly or implicitly.
%In the explicit case, let us consider an orbit containing two nodes in the left hand side of a rule.
%If a node is pulled out to a distinct orbit in the right hand side, then the former orbit
%had split into two dinstinct orbits
For example, the unsewing rule (Fig.~$\ref{fig:rule-unsew-a2}$) splits the edge orbit $L\orbit{\alpha 0, \alpha 2}(n_{0})$ into a pair of edge orbits $R\orbit{\alpha 0, \alpha 2}(n_{0})$ and $R\orbit{\alpha 0, \alpha 2}(n_{1})$.
Therefore, the first application of this rule (Fig.~\ref{fig:unsew-a2-first-unsew}) splits the edge into two edges $H_1\orbit{\alpha 0, \alpha 2}(a_{0})$ and $H_1\orbit{\alpha 0, \alpha 2}(a_{1})$.

%In many cases, a rule rewrites the nodes' labels and, thus, rewrites the way darts
%are connected in a set.
%Such a rewrite may separate them into two or more sets.
A rule can also split an orbit implicitly.
For example, the triangulation rule (Fig.~\ref{fig:rule-triangulation}) splits the face orbit $L\orbit{\alpha 0, \alpha 1}(n_{0})$ along the second implicit links since on the right-hand side, every second implicit link of the nodes  $n_{0}$, $n_{1}$ and $n_{2}$ belonging to the face orbit $R\orbit{\alpha 0, \alpha 1}(n_{0})$ is either deleted or different from $\alpha 0$ and $\alpha 1$.
Therefore, applying this rule (Fig.~\ref{fig:before-after-house-triangulation})
splits $G\orbit{\alpha 0, \alpha 1}(a_{0})$ into four face orbits $H\orbit{\alpha 0, \alpha 1}(a_{0})$,
$H\orbit{\alpha 0, \alpha 1}(c_{0})$, $H\orbit{\alpha 0, \alpha 1}(e_{0})$,
$H\orbit{\alpha 0, \alpha 1}(g_{0})$.


\begin{definition}[Orbit split]\label{def:orbit-split}
  Let $r: L \longrightarrow R$ be a rule, $m: L \to G$ a match,
  $t: G \longrightarrow^{r,m}H$ the transformation of $G$ by $(r,m)$, $\orbit{o}$
  an orbit type, $n$ a preserved node of $r$ and $d$ a preserved dart of $t$.

  An orbit $L\orbit{o}(n)$ is \emph{split} in the rule $r$:
 \begin{itemize}
    \item Either \emph{explicitly} if there exists a preserved node $n'$ in $L\orbit{o}(n)$ that does not belong to the same orbit in $R$, i.e. $R\orbit{o}(n) \neq R\orbit{o}(n')$;

    \item Or \emph{implicitly along the $k$-th implicit link} if there exists in $L\orbit{o}(n)$ a node for which the $k$-th implicit link
          belongs to $\orbit{o}$ and for all nodes $n'$ in $R\orbit{o}(n)$, the $k$-th implicit link is
          renamed outside of $\orbit{o}$.
          % renamed not in $\orbit{o}$.
   \end{itemize}

        An orbit $G\orbit{o}(d)$ is \emph{split} in the transformation $t$ if
        and only if there exists a preserved dart $d'$ in $G\orbit{o}(d)$,
        such that $H\orbit{o}(d) \neq H\orbit{o}(d')$.
\end{definition}

\begin{proposition}\label{prop:orbit-split}
  Let $r : L \longrightarrow R$ be a rule , $m : L \to G$ a match,
  $t : G \longrightarrow^{r,m} H$ the transformation of G by $(r, m)$,
  $\orbit{o}$ an orbit type, $n$ a preserved node of $r$.

  If the orbit $L\orbit{o}(n)$ is both complete and  explicitly split in~$r$, i.e. if there exists a preserved node $n'$ in $L\orbit{o}(n)$ such that $R\orbit{o}(n) \neq R\orbit{o}(n')$, then for any preserved darts $d$ and $d'$ of $t$ respectively matching $R\orbit{o}(n)$ and $R\orbit{o}(n')$,  $G\orbit{o}(d)$ is split in $t$ in two orbits $H\orbit{o}(d) \neq H\orbit{o}(d')$.

  If the orbit $L\orbit{o}(n)$ is either incomplete or implicitly split, the split may or may not appear in $t$ according to $G$.
\end{proposition}

\begin{proofidea}\label{prfid:orbit-split}
  As a reminder, a complete orbit $L\orbit{o}(n)$ entirely matches an orbit
  $G\orbit{o}(d)$.

  In the case of an explicit split in $r$ from $L\orbit{o}(n)$ to $R\orbit{o}(n) \neq R\orbit{o}(n')$, since $L\orbit{o}(n)$ is complete, then due to rules syntactic conditions, $R\orbit{o}(n)$ and $R\orbit{o}(n')$ are also complete.
  Let $d$ and $d'$ be two preserved darts of $t$ such that $d$ is matched by $L\orbit{o}(n)$ and $R\orbit{o}(n)$ and $d'$ is matched by $L\orbit{o}(n)$ and $R\orbit{o}(n')$.
  It follows that $G\orbit{o}(d)$ is split into $H\orbit{o}(d) \neq H\orbit{o}(d')$.

  If $L\orbit{o}(n)$ is complete, and there is an implicit split along the $k$-th implicit links of $L\orbit{o}(n)$ in $r$, then
  by definition, all the $k$-th implicit links of $R\orbit{o}(n)$ are deleted or renamed outside of $\orbit{o}$.
  Let $n'$ be a preserved node of $L\orbit{o}(n)$ with a $k$-th implicit $i$-link with $i$ in $\orbit{o}$, and $d$ and $d'$ two $i$-linked matched darts of $n'$ in $G$.
  Then $d$ and $d'$ may belong to two different orbits in $H$ ($H\orbit{o}(d) \neq H\orbit{o}(d')$); or belong to the same orbit in $H$ if the sub-orbit is folded along the split $i$-link ($H\orbit{o}(d) = H\orbit{o}(d')$); or finally be the same dart if the link is a loop ($d = d'$ and therefore $H\orbit{o}(d) = H\orbit{o}(d')$).

  Let us assume now that $L\orbit{o}(n)$ is not complete.
  Let $d$ and $d'$ be two darts of $G$ matched by $L\orbit{o}(n)$ which are potentially split. %\textit{i.e.} two  darts respectively matching $n'$ and $n''$ in case of explicit splitting of the $i$-link between two preserved nodes $n'$ and $n''$ of $L\orbit{o}(n)$, or two darts matched by a preserved node $n'$ and $i$-linked in $G$ in case of an implicit split along the $k$-th $i$-link of $n'$.
  Because $L\orbit{o}(n)$ is not complete, $G$ can contain some path between $d$ and $d'$ using only links labeled in $\orbit{o}$ such that this path is not entirely matched by $L\orbit{o}(n)$.
  Therefore, this path does not break during the transformation $t$ and both $d$ and $d'$ remain in the same orbit in $H$ ($H\orbit{o}(d') = H\orbit{o}(d'')$).
\end{proofidea}


Referring to the example of the unsewing rule (Fig.~\ref{fig:rule-unsew-a2}), we follow the explicit split of an edge.
The rule splits the edge orbit $L\orbit{\alpha 0, \alpha 2}(n0)$ because both its nodes $n_{1}$ and $n_{2}$ form two distinct edge orbits $R\orbit{\alpha 0, \alpha 2}(n_{0})$ and $R\orbit{\alpha 0, \alpha 2}(n_{1})$.
Rule application (Fig.~\ref{fig:unsew-a2-first-unsew}) entails $n_{0}$ (resp. $n_{1}$) matching both darts $a_{0}$ and $b_{0}$ (resp. $a_{1}$ and $b_{1}$).
The edge orbit $G\orbit{\alpha 0, \alpha 2}(a_{0})$, containing the darts matched by both $n_{0}$ and $n_{1}$, is split into a pair of orbits $H_1\orbit{\alpha 0, \alpha 2}(a_{0})$ (\textit{i.e.} $a_{0}$ and $b_{0}$) and $H_1\orbit{\alpha 0, \alpha 2}(a_{1})$ (\textit{i.e.} $a_{1}$ and $b_{1}$).

Continuing with the face triangulation rule (Fig.~\ref{fig:rule-triangulation}), we follow the implicit split event of a face orbit type $\orbit{\alpha 0, \alpha 1}$.
The orbit $L\orbit{\alpha 0, \alpha 1}(n_{0})$ contains node $n_{0}$ and matches a single set of darts which form a face orbit.
The first implicit link of $n_{0}$ is $\alpha 0$ and the second implicit
link is $\alpha 1$; both belong to orbit type $\orbit{\alpha 0, \alpha 1}$.
However, the orbit $R\orbit{\alpha 0, \alpha 1}(n_{0})$ contains nodes $n_{0}$, $n_{1}$, $n_{2}$ and none of them has either $\alpha 0$ or $\alpha 1$ as a second implicit link. %whose dimension is in $\orbit{\alpha 0, \alpha 1}$.
Under this condition, we can state that the face triangulation rule splits a face into two or more faces.
Applying this triangulation rule on the square face (Fig.~\ref{fig:before-after-house-triangulation}), $L\orbit{\alpha 0, \alpha 1}(n_{0})$ matches the face orbit $G\orbit{\alpha 0, \alpha 1}(a_{0})$).
Then, the square face is split into the four faces matching $R\orbit{\alpha 0, \alpha 1}(n_{0})$:  $H\orbit{\alpha 0, \alpha 1}(a_{0})$, $H\orbit{\alpha 0, \alpha 1}(c_{0})$, $H\orbit{\alpha 0, \alpha 1}(e_{0})$ and $H\orbit{\alpha 0, \alpha 1}(g_{0})$.

Note that both proposition~\ref{prop:orbit-split} and its proof idea mention that a split of an $\orbit{o}$-orbit is guaranteed to happen in a graph transformation when the split is explicit and the orbit is complete in the rule.
Let us consider the unsewing rule (Fig.~\ref{fig:rule-unsew-a2}) and its application on a surface (Fig.~\ref{fig:before-after-unsew-a2}) again but this time, we focus on the surface orbit $L\orbit{\alpha 0, \alpha 1, \alpha 2}(n_{0})$.
The rule nodes of this orbit all miss an $\alpha 1$-link and, thus, $L\orbit{\alpha 0, \alpha 1, \alpha 2}(n_{0})$ is not complete.% $\orbit{\alpha 0, \alpha 1, \alpha 2}$-
Therefore, upon the first application of the rule (Fig.~\ref{fig:unsew-a2-first-unsew}), the surface is not split, although the conditions given in the definition~\ref{def:orbit-split} are met.
Yet, in this case, applying the unsewing rule a second time (Fig.~\ref{fig:unsew-a2-second-unsew}) effectively splits the surface.
Conversely, when an orbit is not complete, the pattern matches only a part of the orbit and statically detected events cannot be guaranteed on $\orbit{o}$.
Therefore, such a statement must be confirmed dynamically.

In the same way, for implicit splits along the $k$-th implicit link, the split is effective only if the matched orbit contains several sub-orbits.
%For example, let be an operation which splits a connected component into volumes. If some connected component is made of only one volume, then the split does not occur.
%of connected component splitting into volumes, the operation does not split in the edge case where this connected component contains only one volume.
%Note that this case of detection, albeit dynamic, is not costly most of the time  since most of practical applications dot not consider orbit folding. Therefore, it is sufficient to check that all the implicit links cut on the rule
%\xavier{pas sûr de comprendre "cut on the rule"}\maxime{dans un cas particulier, les liaisons implicites supposées être scindées sont déjà des boucles, donc il n'y a pas réellement de scission} correspond to loops on the object to transform.
%\agnes{Cela ne devait effectivement pas être clair, car je ne reconnait/comprends plus le §. Je le réécrit donc en français.}
%\agnes{Par exemple, considérons une opération qui déconecte toutes les faces d'une surface. Si elle est appliquée sur une surface qui ne contient qu'une seule face, elle ne provoque aucune scission. Notons que dans ce cas, même si la détection d'une éventuelle scission est dynamique, elle est moins couteuse que dans les méthodes dynamiques, car la scission potentielle est localisée (et donc la vérification à faire plus précise). De plus, dans de nombreuses applications, les orbites ne peuvent pas être repliées sur elle-mêmes. Une orbite est alors $i$-isolées si toutes ses $i$-liaisons sont des boucles. Dans notre exemple, si la surface ne contient aucune arrête donc les deux cotés sont incidents à la même face, alors la surface ne contient qu'une seule face si et seulement si toutes ses $2$-liaisons sont des boulces.}
For instance, let us consider an operation which disconnects every face from a shell surface. Applying such an operation onto a single face has no effect: this face is not split, but a dynamic verification must nevertheless be carried out afterwards. Fortunately, even if detecting a potential split is a dynamic process, it is less time-consuming than classical methods, since we know exactly upon which cell the split could occur.
%\agnes{C'est même plus précis que cela on sait entre quel couple de brins la scission peut apparaitre.}

%Moreover, in most practical applications, orbits are not folded on themselves. An orbit is "$i$-isolated" if all its $i$-links are loops. Using the example above, if the shell surface does not contain any edge where both sides are adjacent to the same face, then the shell contains only one face if and only if each of its links are loops.

%\agnes{Cela reste tout de même complexe, donc on peut supprimer la dernière partie (pour faire vite) ou ajouter des dessins (si on a le temps).}
%\xavier{oui, sans dessin c'est compliqué. On peut alors arrêter le texte à "since we know exactly upon which cell the split could occur".}
%\agnes{Oui c'est l'idée.}

\subsection{Merging}

\begin{figure}[t]
  \centering
  \maxsizebox{\columnwidth}{.15\columnwidth}{\includegraphics{rule-sewA2.pdf}}
  \caption{Rule $r : L \longrightarrow R$ of $\alpha 2$-sewing along two edges
    ($\orbit{\alpha 0, \alpha 2}$)-orbit}
  \label{fig:rule-sewing-a2}
\end{figure}

\begin{figure}[t]
    \centering
    %\includegraphics[width=.65\columnwidth]{rule-delete-edge.pdf}
    \maxsizebox{\columnwidth}{.15\columnwidth}{\includegraphics{rule-delete-edge.pdf}}
    \caption{Rule $r: L \longrightarrow R$ of edge deletion}
    \label{fig:rule-delete-edge}
\end{figure}

\begin{figure}[t]
  \begin{subfigure}[t]{.45\columnwidth}
    \centering
    \includegraphics[height=.25\textheight]{demo-fig-delete-edge-left-match.pdf}
    \caption{$G$}
    \label{fig:house-merge-left-match}
  \end{subfigure}
  \hfill
  \begin{subfigure}[t]{.45\columnwidth}
    \centering
    \includegraphics[height=.25\textheight]{demo-fig-delete-edge-right-match.pdf}
    \caption{$H$}
    \label{fig:house-merge-right-match}
  \end{subfigure}
  \caption{Application of the edge deletion rule \ref{fig:rule-delete-edge} on match dart $a_{0}$}
  \label{fig:before-after-delete-edge}
\end{figure}

Based on the same intuition as for orbit split, there exist explicit and implicit merging orbits.

For example, in the explicit case, the sewing rule (Fig.~\ref{fig:rule-sewing-a2}) is the opposite of the unsewing rule (Fig.~\ref{fig:rule-unsew-a2}) in regard of both the result of its application and its construction.
Let us consider the square plan (Fig.~\ref{fig:before-after-unsew-a2}) from $H_{1}$ to $G$ (Fig.~\ref{fig:unsew-a2-first-unsew} to Fig.~\ref{fig:unsew-a2-left-match}).
The sewing rule merges two edge orbits $L\orbit{\alpha 0, \alpha 2}(n_{0})$ and $L\orbit{\alpha 0, \alpha 2}(n_{1})$ into a single edge orbit $R\orbit{\alpha 0, \alpha 2}(n_0)$.
Therefore,  applying  the sewing rule merges two edges of $H_{1}$, $H_1\orbit{\alpha 0, \alpha 2}(a_0)$ and $H_1\orbit{\alpha 0, \alpha 2}(c_1)$, into one edge $G\orbit{\alpha 0, \alpha 2}(a_{0})$.

The same goes for the implicit merging.
For example, in the implicit case, the edge deletion rule (Fig.~\ref{fig:rule-delete-edge}) merges the face vertex orbits mapped to $L\orbit{\alpha 1}(n_{0})$ along the second implicit link.
In fact, in the left-hand side of the rule, the second implicit link of $n_0$ is $\alpha 2$ % out of orbit $\orbit{\alpha 1}$
and $n_1$ has no second implicit link.
In the right-hand side of the rule, the second implicit link of $n_{1}$ is rewritten as $\alpha 1$.
Therefore, the application of this rule (Fig.~\ref{fig:before-after-delete-edge}) merges $G\orbit{\alpha 1}(a_{1})$ with $G\orbit{\alpha 1}(b_{1})$ into $H\orbit{\alpha 1}(a_{1})$ and $G\orbit{\alpha 1}(d_{1})$ with $G\orbit{\alpha 1}(c_{1})$ into $H\orbit{\alpha 1}(d_{1})$.


\begin{definition}[Orbit merging]
    Let $r : L \longrightarrow R$ be a rule, $m : L \to G$ a match, $t : G \longrightarrow^{r,m} H$ a transformation of $G$ by $(r,m)$, $\orbit{o}$ an orbit type, $n$ a preserved node of $r$ and $d$ a preserved dart of $t$.

    An orbit $L\orbit{o}(n)$ is \emph{merged} by the rule $r$:
    \begin{itemize}
        \item Either \emph{explicitly} if there exists a preserved node $n'$, with $L\orbit{o}(n) \neq L\orbit{o}(n')$, such that both $n$ and $n'$ belong to the same orbit in $R$, i.e. $R\orbit{o}(n) = R\orbit{o}(n')$.

        \item Or \emph{implicitly} along the $k$-th implicit link if there exists in $R\orbit{o}(n)$ a node whose $k$-th implicit link belongs to $\orbit{o}$ and for all nodes $n'$ in $L\orbit{o}(n)$, the $k$-th implicit link does not belong to $\orbit{o}$.
    \end{itemize}

    Two different orbits $G\orbit{o}(d)$ and $G\orbit{o}(d')$, with $d'$ a preserved dart, are \emph{merged} in the transformation $t$ if and only if $H\orbit{o}(d) = H\orbit{o}(d')$.
\end{definition}

\begin{proposition}% \label{prop:orbit-merge}
   Let $r : L \longrightarrow R$ be a rule, $m : L \to G$ a match,
  $t : G \longrightarrow^{r,m} H$ a transformation of $G$ by $(r,m)$,
  $\orbit{o}$ an orbit type, $n$ a preserved node of $r$.

  If the orbit $L\orbit{o}(n)$ is both complete and explicitly merged, \textit{i.e.} if there exists a preserved node $n'$ in $r$ such that $L\orbit{o}(n) \neq L\orbit{o}(n')$ and $R\orbit{o}(n) = R\orbit{o}(n')$, then for any preserved darts $d$ and $d'$ of $t$ respectively matched by $L\orbit{o}(n)$ and $L\orbit{o}(n')$, $G\orbit{o}(d)$ is merged with $G\orbit{o}(d')$ in $H$, \textit{i.e.} $G\orbit{o}(d) \neq G\orbit{o}(d')$ and $H\orbit{o}(d) = H\orbit{o}(d')$.

  If the orbit $L\orbit{o}(n)$ is either not complete or implicitly merged, the merge may or may not appear in $t$ according to~$G$.
\end{proposition}

\begin{proofidea}
  The idea is analogous to the one developed for Proposition~\ref{prop:orbit-split}.

  If the orbit $L\orbit{o}(n)$ is both complete and explicitly merged with another orbit $L\orbit{o}(n')$, where $n'$ is a preserved node of $r$, then $R\orbit{o}(n) = R\orbit{o}(n')$.
  Because $L\orbit{o}(n)$ is complete, $L\orbit{o}(n)$ entirely matches $G\orbit{o}(d)$ for any preserved dart $d$ of $t$ matched by $L\orbit{o}(n)$.
  Due to the injection property of matching, for any preserved dart $d'$ of $t$ matched by $L\orbit{o}(n')$, $d'$ is not a dart of $G\orbit{o}(n)$.
  In other words, $G\orbit{o}(n) \neq G\orbit{o}(n')$.
  Finally, due to rule application,  $R\orbit{o}(n) = R\orbit{o}(n')$ implies $H\orbit{o}(d) = H\orbit{o}(d')$.

  If $L\orbit{o}(n)$ is complete and if there is an implicit merge along the $k$-th implicit links of $L\orbit{o}(n)$ in $r$, then by definition, all the $k$-th implicit links of $L\orbit{o}(n)$ are deleted or named outside of~$\orbit{o}$.
  Let $n'$ be a preserved node of $L\orbit{o}(n)$ such that the $k$-th implicit link of $n'$ is $i$-labeled in $R$ with $i$ in $\orbit{o}$.
  Let $d$ and $d'$ be two darts, matched by $n'$, and $i$-linked in $H$.
  $d$ and $d'$ belong to the same $\orbit{o}$-orbit in $H$ ($H\orbit{o}(d) = H\orbit{o}(d')$).
  Then $d$ and $d'$ may belong to two different orbits in $G$ ($G\orbit{o}(d) \neq G\orbit{o}(d')$); or belong to the same orbit in $G$ if the rule folds the sub-orbit along the $i$-link ($G\orbit{o}(d) = G\orbit{o}(d')$); or finally be the same dart if the $i$-link is a loop ($d = d'$ and therefore  $G\orbit{o}(d) = G\orbit{o}(d')$).

  Let us assume now that $L\orbit{o}(n)$ is not complete. Let $d$ and $d'$ be two preserved darts of $t$ matched by $H\orbit{o}(d)$ which are potentially merged.
  Because $L\orbit{o}(n)$ is not complete, $R\orbit{o}(n)$ is also not complete, $H$ can contain some path between $d$ and $d'$ using only links in $\orbit{o}$ such that this path is not entirely matched by $R\orbit{o}(n)$.
  Therefore, this path has not been built during the transformation $t$ and $d$ and $d'$ were already in the same orbit in $G$ ($G\orbit{o}(d) = G\orbit{o}(d')$).
\end{proofidea}

With the example of the sewing rule (Fig.~\ref{fig:rule-sewing-a2}), we follow the explicit merging of two edges.
The edge orbits $L\orbit{\alpha 0, \alpha 2}(n_{0})$ and $L\orbit{\alpha 0, \alpha 2}(n_{1})$ are complete and merged in the rule because their respective nodes $n_{1}$ and $n_{2}$ are part of the same edge orbit $R\orbit{\alpha 0, \alpha 2}(n_{0})$.
Applying the transformation, $n_{0}$ (resp. $n_{1}$) matches the darts $a_{0}$ and $b_{0}$ (resp. $a_{1}$ and $b_{1}$).
Edge orbits $G\orbit{\alpha 0, \alpha 2}(a_{0})$ and $G\orbit{\alpha 0, \alpha 2}(a_{1})$ are merged into orbit $H\orbit{\alpha 0, \alpha 2}(a_{0})$ (\textit{i.e.} $a_{0}$, $b_{0}$, $a_{1}$ and $b1_{0}$), as shown in Fig.~\ref{fig:unsew-a2-left-match}.

With the edge deletion rule (Fig.~\ref{fig:rule-delete-edge}), we follow the implicit merging event of the face vertices of orbit types $\orbit{\alpha 1}$.
The orbit $L\orbit{\alpha 1}(n_{0})$ contains both nodes $n_{0}$ and $n_{1}$ and matches a single set of darts representing four face vertices.
In this example, since all nodes of $L\orbit{\alpha 1}(n_{0})$ miss an implicit $\alpha 1$ link and since the only remaining node in $R\orbit{\alpha 1}(n_{1})$ has such a link, then we can state that face vertices are merged.
Applying this sewing rule on the horizontal edge of the house (Fig.~\ref{fig:before-after-delete-edge}), $L\orbit{\alpha 1}(n_{0})$ matches face vertices $G\orbit{\alpha 1}(a_{1})$, $G\orbit{\alpha 1}(b_{1})$, $G\orbit{\alpha 1}(c_{1})$ and $G\orbit{\alpha 1}(d_{1})$.
Then, the face vertices around the edge to delete will be merged into the pair of face vertices matched from $R\orbit{\alpha 1}(n_{1})$ (\textit{i.e.} $H\orbit{\alpha 1}(a_{1})$ and $H\orbit{\alpha 1}(d_{1})$).

Note that although the edge deletion rule may merge two adjacent faces (Fig.~\ref{fig:house-merge-right-match}), the rule is not complete on $\orbit{\alpha 0, \alpha 1}$ since the first implicit link of $n_1$ in $R$ does not belong to $\orbit{o}$.
Therefore, such a merge must be confirmed dynamically.


\subsection{Other events}
\label{sec:other-events}

The remaining events, namely \emph{deletion}, \emph{non modification} and
\emph{modification}, will be presented in a concise way.

\subsubsection{Deletion}

This event mirrors the creation one. Let us get back to the edge deletion rule (Fig.~\ref{fig:rule-delete-edge}).
This rule, as its name suggests, deletes an edge orbit $L\orbit{\alpha 0, \alpha 2}(n_{0})$.
Its single node $n_{0}$, which by itself matches an edge because it is complete on $\orbit{\alpha 0, \alpha 2}$, no longer exists on the right side of the rule and, therefore, is deleted.
Applying the rule on the house (Fig.~\ref{fig:before-after-delete-edge}), $L\orbit{\alpha 0, \alpha 2}(n_{0})$ matches the orbit $G\orbit{\alpha 0, \alpha 2}(a_{0})$ whose darts ($a_0$ $b_0$, $c_0$ and $d_0$) are then deleted in $H$ (Fig.~\ref{fig:house-merge-right-match}), since there are no nodes in $R$ to match them.

\begin{definition}[Orbit deletion]
    Let $r : L \longrightarrow R$ be a rule, $m : L \to G$ a match, $t : G \longrightarrow^{r,m} H$ the transformation of $G$ by $(r,m)$, $\orbit{o}$ an orbit type, $n$ a node of $L$ and $d$ a dart of $G$.

    An orbit $L\orbit{o}(n)$ is \emph{deleted} in the rule $r$ if all of its nodes are deleted in $r$.

    An orbit $G\orbit{o}(d)$ is \emph{deleted} in the transformation $t$ if and only if all of its darts are deleted by $t$.
\end{definition}

\begin{proposition}
   Let $r : L \longrightarrow R$ be a rule, $m : L \to G$ a match, $t : G \longrightarrow^{r,m} H$ the transformation of $G$ by $(r,m)$, $\orbit{o}$ an orbit type, $n$ a node of $L$ and $d$ a dart of $G$.

   An orbit $L\orbit{o}(n)$ is deleted in the rule $r$ if and only if the orbit $G\orbit{o}(d)$ is deleted in the transformation $t$ for any dart $d$ matched by $L\orbit{o}(n)$.
\end{proposition}

\begin{proofidea}
  Thanks to the syntactic conditions on the rules, a deleted orbit is always complete.
  Consequently, the corresponding orbits are also completely deleted in the transformation.
\end{proofidea}


\subsubsection{Non-Modification}

An orbit is said to be \emph{not modified} when its nodes and links remain unchanged.
For example, in the triangulation rule (Fig.~\ref{fig:rule-triangulation}), the face edge orbit $L\orbit{\alpha 0}(n_{0})$ only contains node $n_{0}$  and the same goes for the orbit $R\orbit{\alpha 0}(n_{0})$.
Considering that the implicit $\alpha 0$ link of $n_{0}$ is unchanged and that there is no
other node in its orbit, we can state that the triangulation rule does not modify the face edges.
Applying this rule on the house (Fig.~\ref{fig:before-after-house-triangulation}), the face edge orbits remain the same and $G\orbit{\alpha 0}(a_{0}) = H\orbit{\alpha 0}(a_{0})$, $G\orbit{\alpha 0}(c_{0}) = H\orbit{\alpha 0}(c_{0})$, $G\orbit{e_{0}} = H\orbit{e_{0}}$ and $G\orbit{\alpha}(g_{0}) = H\orbit{\alpha 0}(g_{0})$ (Fig.~\ref{fig:house-triangulation-right-match}).

\begin{definition}[Orbit non-modification]
    Let $r : L \longrightarrow R$ be a rule, $m : L \to G$ a match, $t : G \longrightarrow^{r,m} H$ the transformation of $G$ by $(r,m)$, $\orbit{o}$ an orbit type, $n$ a preserved node of $r$ and $d$ a preserved dart of $t$.

    An orbit $L\orbit{o}(n)$ remains \emph{not modified} in the rule $r$ if $L\orbit{o}(n) = R\orbit{o}(n)$ and all the nodes of the orbit have the same label on both sides of the rule, \textit{i.e.} their implicit links are the same in $L$ and $R$.

    An orbit $G\orbit{o}(d)$ is said to be \emph{not modified} in the transformation $t$ if $G\orbit{o}(d) = H\orbit{o}(d)$.
\end{definition}

\begin{proposition}
    Let $r : L \longrightarrow R$ be a rule, $m : L \to G$ a match, $t : G \longrightarrow^{r,m} H$ the transformation of $G$ by $(r,m)$, $\orbit{o}$ an orbit type, $n$ a preserved node of $r$ and $d$ a preserved dart of $t$.

    If an orbit $L\orbit{o}(n)$ remains not modified in the rule $r$, then $G\orbit{o}(d)$ remains not modified in the transformation $t$ for any dart $d$ of $G$ matched by $L\orbit{o}(n)$.

    If an orbit $G\orbit{o}(d)$ is not matched by the rule, \textit{i.e.} no dart of $G\orbit{o}(d)$ is matched by $L$, then it remains not modified
\end{proposition}

\begin{proofidea}
   This is obvious, because a rule may modify only the matched part of a graph.
\end{proofidea}


\subsubsection{Modification}

%This event mirrors the one about permanent.
Generally speaking, an orbit is \emph{modified} when it is neither created, deleted, not modified, split or merged.
More specifically, a modified orbit has either modified links, some created part or some deleted part.
Also, such a change must not lead to a split nor a merging event.
For example, once again considering the triangulation rule,  vertex orbit $L\orbit{\alpha 1, \alpha 2}(n_{0})$ only contains  node $n_{0}$, while $R\orbit{\alpha 1, \alpha 2}(n_{0})$ contains both nodes $n_{0}$ and $n_{1}$.
The orbit has a created part related to $n_{1}$, the second implicit link of $n_{0}$ is deleted, and there is no split, either implicit or explicit: thus, the orbit is modified.
Applying this rule on the house (Fig.~\ref{fig:before-after-house-triangulation}) entails vertex orbits $G\orbit{\alpha 1, \alpha 2}(a_{0})$, $G\orbit{\alpha 1, \alpha 2}(c_{0})$, $G\orbit{\alpha 1, \alpha 2}(e_{0})$, $G\orbit{\alpha 1, \alpha 2}(g_{0})$ to be modified with an added part of two new darts each.

%\begin{definition}[Orbit modification]
%    Let $r : L \longrightarrow R$ be a rule, $m : L \to G$ a match, $t : G \longrightarrow^{r,m} H$ the transformation of $G$ by $(r,m)$, $\orbit{o}$ an orbit type, $n$ a preserved node and $d$ a preserved dart matched by $n$.

%    An orbit $L\orbit{o}(n)$ is modified in the rule $r$ if its reduction/augmentation or its nodes relabeling does not produce any of the other events in $r$.

%    An orbit $G\orbit{o}(d)$ is modified in the transformation $t$ if its only augmented or reduced by $t$.
%\end{definition}

%%% Local Variables:
%%% mode: latex
%%% TeX-master: "../main"
%%% End:
